\section{模型评价}

\subsection{模型优点}

\textbf{(1) 结构统一、约束对齐题目。} 四个问题都按“预测 $\rightarrow$ 优化 $\rightarrow$ 因果回放”推进：问题一是确定性线性规划，问题二到四是滚动线性规划加锁死合同的因果回放。目标函数统一为购电费，约束与题目逐条对应——储电量运行窗 $1200$--$10800$ kWh、充放电效率 $0.9$、接口功率 $5000$ kW、紧急购电 $5$ 倍、调整部分少买 $0.5$ 倍/多买 $1.5$ 倍。任一时刻的购电、充放电、弃光和紧急购电都能追溯到具体约束。

\textbf{(2) 可复现。} 算法与超参数全部固定（XGBoost 随机种子 $42$、线性规划均求到最优、参数冻结），并配有泄漏检查（预测只用 $D$ 日前信息，错误 $0$ 条）、价格泄漏检查（$0$ 条）与逐时段约束校验（能量平衡、SOC、功率、同时充放电均为 $0$ 错误）。同一套代码重复运行得到同样数字。

\textbf{(3) 有对照、参数在样本外前标定。} 每个问题都设了可对照基线：问题一对照无储能与分时贪婪；问题二对照无分位、固定分量分位，再在净分位上叠加日末约束与保护回放；问题三对照不调整的 N0 与选择门 M0L；问题四对照固定裕度与完美电价 oracle。问题二的 $\alpha$、$\rho$、$S^{*}$ 在 $1$ 月 $15$--$31$ 日网格选定，问题三的 $q_L$ 与展望长度在标定区间选定后再做全年回测。

\subsection{模型不足与推广}

\textbf{(1) 预测依赖历史平稳性。} 负载、光伏与电价都用“历史同刻均值 + XGBoost”类统计模型，特征只有时间，不含天气、节假日与极端事件。样本外出现未见过形态时，偏差会被 $5$ 倍紧急电价放大；结论基于附件 2/附件 4 的 2025 年样本，换年份需重新标定。

\textbf{(2) 若干简化假设。} 充放电效率恒为 $0.9$、储能无自放电与寿命损耗、电价与负载相互独立、$10$ min 内功率恒定；寿命成本未计入目标，实际费用可能因此偏高或偏低。

\textbf{(3) 回放是因果近似而非全知最优。} 日内回放用贪心与滚动 LP，是在“不利用未来信息”下的近似；问题四完美电价诊断显示，距用真实电价决策还差约 $0.57\%$，实际执行偏差也会带来额外费用。

\textbf{(4) 参数在标定区间内选定。} 问题二的 $S^{*}$ 与 $\rho$ 只在同年 $1$ 月 $16$ 格选定，问题三的分位与缓冲系数在各自区间内选定，并非独立年份的样本外检验。可改进方向：给预测加天气特征、用随机规划或鲁棒优化刻画尾部风险、把储能损耗计入目标，并做多年份滚动验证。